% Compile with:
%   pdflatex paper_alife && bibtex paper_alife && pdflatex paper_alife && pdflatex paper_alife
\documentclass[11pt]{article}

\usepackage[a4paper, margin=1in]{geometry}
\usepackage{amsmath,amssymb}
\usepackage{graphicx}
\usepackage{booktabs}
\usepackage{caption}
\usepackage{subcaption}
\usepackage{natbib}
\usepackage{xcolor}
\usepackage[hidelinks]{hyperref}
\usepackage[final,protrusion=true,expansion=false]{microtype}

\graphicspath{%
  {../../out/figures/phase/}%
  {../../out/figures/scenarios/}%
  {../../out/figures/sweeps/}%
  {../../out/figures/policies/}%
}

\title{Productive-flow share and wealth-coupled cultural transmission:\\
\large a stylised agent-based model of cultural erosion under economic pressure}

\author{[Author]\\
\small (supervision and acknowledgements TBD)}
\date{Draft of \today\ for \emph{Artificial Life} / ALIFE}

\begin{document}
\maketitle

\begin{abstract}
We present an agent-based model of an economy with seven occupationally
heterogeneous agent classes that exchanges its monetary state through
production, extraction, credit, taxation, inheritance, and a
bandwidth-limited cultural-transmission channel. The cultural channel
couples to wealth on a log scale: the number of skill dimensions a child
inherits from a parent rises as the log of parent wealth, so the
transmission bandwidth is graded across orders of magnitude of wealth
rather than saturating at a single threshold.

We argue and demonstrate three results that we believe make the model
useful for an artificial-life audience.

First, the model is honest about its scope. The K=5 skill space is
\emph{closed} by construction: mutation can reset a dimension to a fresh
draw but no genuinely new skill dimension can appear. The model
therefore exhibits cultural \emph{erosion} under wealth concentration
rather than cumulative cultural evolution. This is the right framing
relative to \citet{henrich2015secret} and \citet{henrich2004demography},
with \citet{vaesen2016cumulative} noted as a counter-case. We present a
minimal open-ended extension that allows new skill dimensions to enter
the global pool via wealth-biased novelty events, and show on a
per-capita basis that the discovery channel concentrates novelty in
wealthier lineages.

Second, the cultural-economic asymmetry of the model, in which cultural
transmission alone does not equalise wealth even when equalised at
birth, is grounded in mechanism rather than in architectural wiring.
A class-by-class regression of 40-tick wealth gain on the skill vector
shows that workers and makers have strong skill-to-wealth fits ($R^2
\geq 0.97$ with normalised elasticities $\eta \geq 0.76$ on the
class-typical dimensions); rentiers, lenders, and speculators have
near-zero fits ($R^2 \leq 0.13$, $|\eta| \leq 0.10$). The
cultural-economic asymmetry is therefore the consequence of one half of
the economy (extractive classes) having weak internal skill-wealth
feedback while the other half (productive classes) has strong feedback.

Third, we introduce a continuous per-tick scalar, the
\emph{productive-flow share} $\Phi_t = F^\text{prod}_t /
(F^\text{prod}_t + F^\text{ext}_t)$, which captures whether per-tick
monetary accrual is dominated by productive- or extractive-class activity.
$\Phi$ is not an autocatalytic-set closure in the
\citet{kauffman1988evolution} / \citet{hordijk2004detecting} RAF sense;
it is a flow-share measure operationalising the productive-vs-extractive
distinction that the heterodox economic tradition has long named
without instrumenting. The Sankey of steady-state per-tick flows
makes $\Phi$ legible as a diagram rather than as a scalar.

We note in passing a result from the policy-comparison side of the
broader project: an absolute cap on inheritance raises the wealth
share held by the bottom 50\% of the population from 0.0\% under
status quo to 19.2\% (cap alone) or 25.3\% (cap plus 2\% wealth
tax).  The same intervention raises the productive-flow share
$\Phi$ from approximately zero to 0.27--0.37.  Among the
instruments we tested, the cap is the only one that changes both
the ownership distribution and the productive--extractive flow
balance materially.  The cultural-economic asymmetry result of
this paper provides part of the mechanistic explanation: the
extractive classes have weak skill-to-wealth feedback, so
redistributing skills does not redistribute wealth, but
redistributing the inheritance of wealth itself does.  The
policy-discourse implications of the cap result are reported in
the companion submission targeting Fiscal Studies.

We position these cultural contributions relative to the
cultural-evolution and
open-ended-evolution literatures, locate the model relative to
\citet{spangler2022distribution} as the closest published economic-ABM
neighbour, and discuss the limits of the cultural model and the
open-skills extension. The fiscal-policy results that follow naturally
from the model are summarised in a companion submission targeting
the heterodox-economics literature.
\end{abstract}

%==========================================================================
\section{Introduction}
%==========================================================================

The artificial-life community has long studied open-ended evolution,
cumulative cultural transmission, and the autocatalytic structure of
chemical and biological systems. This paper presents a model that
borrows from each of those traditions but applies them to a question
the alife literature has not previously taken on directly: how
economic inequality and cultural diversity co-evolve, and what
mechanisms cause the cultural sphere to track the economic sphere.

We approach this with a stock-flow agent-based simulation. Seven
occupationally typed agents (workers, makers, employers, lenders,
rentiers, speculators, and a single state) exchange a monetary state
variable through production, extraction (rent, interest, capital
return), inheritance, and a bandwidth-limited cultural-transmission
channel. Each agent carries a five-dimensional skill vector; on death
a same-class child inherits a number of skill dimensions that scales
with the log of parent wealth. The mechanism is simple, and it has the
property that cultural inheritance concentrates in wealthier lineages
not by an explicit class bias but by the wealth-coupling of the
bandwidth gate itself.

We organise the paper around three contributions.

The first is honesty about scope. The K=5 skill space is closed by
construction. In this respect the model is squarely on one side of the
open-endedness debate
\citep{taylor2019evolutionary, banzhaf2016defining}: the only thing that
can evolve in the main model is the wealth distribution and the subset
of the existing skill pool held by each agent. The
\citet{henrich2015secret} ``cumulative culture'' framing therefore
does not strictly apply. What the model demonstrates is the
\emph{erosion} of cultural breadth under wealth concentration, in the
spirit of \citet{henrich2004demography}, with
\citet{vaesen2016cumulative}'s critique of the Tasmanian inference
noted. We add (\S\ref{sec:open-skills}) a stand-alone open-ended
extension in which new skill dimensions can enter the pool via
wealth-biased novelty events.

The second is the mechanism behind the cultural-economic asymmetry.
The model exhibits a striking negative result: cultural transmission
alone does not equalise wealth, even when transmission is equalised at
birth. An obvious worry is that this asymmetry is built into the
architecture. We address that worry by quantifying the
skill-to-wealth channel directly with a class-by-class regression of
wealth gain on the skill vector. The asymmetry turns out to be partly
structural: the extractive classes (rentiers, lenders, speculators)
have $R^2$ between 0.11 and 0.13 and normalised elasticities of order
$0.05$ to $0.10$. But the productive classes (workers, makers) have
$R^2$ of $0.97$ to $0.98$ and normalised elasticities of $0.76$ to
$0.95$ on the class-typical dimensions. The asymmetry is therefore
not an arbitrary wiring choice but the consequence of one half of the
economy having strong internal skill-to-wealth feedback and the other
half not. The model is doing the right thing.

The third is a measurement device. We define $\Phi_t$, the
productive-flow share, as the share of per-tick monetary accrual
accruing to or being created by productive-class agents:
$\Phi_t = F^\text{prod}_t / (F^\text{prod}_t + F^\text{ext}_t)$.
We use this to characterise regimes (Section~\ref{sec:phase}) and to
visualise wealth flows (the Sankey, Section~\ref{sec:sankey}). We
explicitly distinguish $\Phi$ from the topological autocatalytic-set
closure of \citet{kauffman1988evolution} and
\citet{hordijk2004detecting}, after which it is named in motivation but
not in technical content.

The fiscal-policy implications of the model (wealth-tax design,
inheritance trapdoor, calibrated-vs-default policy comparisons) are
substantial and are reported separately in a companion submission;
this paper engages them only to the extent needed to give a cultural-
evolutionary reading of the results.

%==========================================================================
\section{Related work}
\label{sec:related}
%==========================================================================

The most direct alife-adjacent precedent on cultural transmission is
the cultural-evolution tradition of
\citet{cavallisforza1981cultural}, \citet{boyd1985culture},
\citet{henrich2015secret}, and \citet{lewis2012transmission}.
The wealth-coupled bandwidth mechanism we introduce sits adjacent to:
\citet{henrich2001evolution} on prestige-biased transmission (where
the \emph{target} of imitation is status-biased; we couple
the \emph{bandwidth} of transmission to wealth instead);
\citet{henrich2004demography} on population-size-capped cumulative
culture, where \citet{vaesen2016cumulative} provide the counter-
analysis arguing the relevant population sizes in the Tasmania case
were not as different as the model required;
\citet{lewis2012transmission} on the empirical sensitivity of
cumulative culture to transmission fidelity; and
\citet{charbonneau2021fidelity} on the ``grain problem'' for fidelity
measures, which directly applies to our K=5 discretisation.
We have not been able to locate a published agent-based model that
couples transmission \emph{bandwidth} to wealth specifically.

On the alife/OEE side we engage with \citet{hashimoto2019open} on
open-ended evolution of tag combinations,
\citet{taylor2019evolutionary} on routes to open-ended evolution, and
\citet{banzhaf2016defining} on the OEE program more generally. The
TAP-style adjacent-possible framing of
\citet{hordijk2023emergence} motivates our open-ended skill-discovery
extension.

The closest economic-ABM neighbour is \citet{spangler2022distribution},
who couple estate taxation with binary skill inheritance and a Carnegie
effect. We test their headline mechanism directly under a clean
toggle in \S\ref{sec:spangler}.

For the productive-vs-extractive framing of $\Phi$ we draw on
\citet{kauffman1988evolution} and \citet{mazzucato2018value};
\citet{hordijk2023emergence} is the closest neighbour on detecting
autocatalytic structure in production-function networks, where they
use a topological RAF detection and we use a continuous flow ratio.

%==========================================================================
\section{The model}
\label{sec:model}
%==========================================================================

\subsection{Agent types and population}

Seven agent classes are present at fixed proportions per 100 agents:
72 workers, 9 makers, 7 employers, 5 lenders, 3 rentiers, 3
speculators, and 1 state. Each agent carries wealth, debt, age,
lifespan drawn from $\mathcal{N}(70, 12)$ truncated to $[20, 120]$,
a five-dimensional skill vector with class-typical biases at
initialisation, peak lifetime wealth, and a bandwidth-at-birth record
used for the analysis in \S\ref{sec:bandwidth-decomp}.

The skill vector dimensions are interpreted as \emph{technical},
\emph{organisational}, \emph{social}, \emph{financial}, and
\emph{generative}. Initialisation gives each class a class-typical
prior bias on a subset of dimensions; the bias is small and is
preserved by the inheritance rule below.

The model runs at $N = 100$ in the cultural results reported here.
A scale-up to $N = 10{,}000$ is available; the cultural results are
qualitatively unchanged under scaling and we report them at $N=100$
for legibility.

\subsection{Per-tick mechanics, in brief}

Each tick runs the following phases in order. Production: makers and
employers hire workers from a skill-sorted queue; makers transform
a material input plus paid labour into output at multiplier $M$;
employers capture a markup. State optionally hires unemployed
workers and purchases maker output. Extraction: workers pay rent
to rentiers in proportion to rentier wealth; rentiers earn capital
return $r$; speculators accrue a stochastic intra-stock return.
Consumption and credit: subsistence is debited from every non-state
agent; shortfalls are borrowed from a lender pool, with starvation if
the pool is exhausted. Interest and repayment: debtors pay interest
when liquid, principal when wealth allows. Tax and UBI. Life: agents
whose age exceeds their lifespan die naturally; a child of the same
class is born with inherited wealth (after inheritance tax and cap)
and bandwidth-limited inherited skills (next subsection).
A complete ODD-style summary is in Appendix~\ref{app:odd}.

\subsection{Cultural transmission: bandwidth-limited inheritance}
\label{sec:transmission}

When a parent agent of class $t$ dies and is replaced by a
same-class child, the child's skill vector is constructed as follows.
The bandwidth, i.e.\ the number of skill dimensions the child
inherits from the parent, is
\begin{equation}
\text{bw\_dims}(w_{\text{parent}}) \;=\; \text{clip}\!\Big(\,
b_0 K \;+\; b_w K \cdot \log_{10}\!\big(\max(1, w_{\text{parent}})\big) / 4,
\;\; 1, \; K \,\Big),
\label{eq:bandwidth}
\end{equation}
with $K=5$, $b_0$ a base proportion ($b_0 = 0.5$ default), and $b_w$
a wealth-coupling slope ($b_w = 0.5$ default). The child inherits
the parent's $\text{bw\_dims}$ strongest skill dimensions; the rest
are zero. A small per-dimension mutation rate ($\mu = 0.05$ default)
replaces an inherited value with a fresh Beta(2,2) draw on its own
account.

The log-wealth coupling is the key mechanism for this paper. With
parent wealth at $1$ unit, $\log_{10}(1) = 0$, so $\text{bw\_dims}
= b_0 K = 2.5$ (rounded to 3). At parent wealth $10{,}000$ units,
$\log_{10}(10000)/4 = 1$, so $\text{bw\_dims} = b_0 K + b_w K = 5$
(all skills transmitted). The coupling is thus graded across four
orders of magnitude of parent wealth, rather than saturating at a
single intermediate threshold as a prior linear formulation did.

\emph{Knowledge loss} is recorded as the per-tick sum of skill mass
present in each dying parent but not transmitted to its child.

\subsection{The productive-flow share metric}
\label{sec:phi}

During each tick we accumulate two flow sums treated symmetrically.
Productive accrual $F^\text{prod}_t$ sums (i) wages paid by makers,
employers, and the state to workers; (ii) state purchases of maker
output; (iii) maker value-add, i.e. the surplus $\text{output} -
\text{inputs}$ created in production and retained by the maker;
(iv) a fraction $0.4$ of gross employer markup, treated as productive
value-add reflecting the employer's organising work. Extractive
accrual $F^\text{ext}_t$ sums rent paid by workers to rentiers,
interest paid by debtors to lenders, capital returns accruing to
rentier stocks, and positive speculator returns (negatives not
counted, by the convention by which rent and interest do not have
negative counterparts).

The productive-flow share is
\begin{equation}
\Phi_t \;=\; \frac{F^\text{prod}_t}{F^\text{prod}_t + F^\text{ext}_t} \in [0, 1].
\label{eq:phi}
\end{equation}

\paragraph{Nomenclature.} $\Phi$ is \emph{not} an autocatalytic-set
closure in the technical sense of
\citet{kauffman1988evolution} or \citet{hordijk2004detecting}, who
detect a topological RAF property of a chemical or production graph.
$\Phi$ is a continuous, scalar, per-tick flow-share measure
operationalising the productive-vs-extractive distinction. We have
named the metric ``productive-flow share'' to avoid friction with
the RAF-literate audience.

%==========================================================================
\section{The phase portrait and the regime classification}
\label{sec:phase}
%==========================================================================

We characterise the model by running thirteen named scenarios at
$N=100$ for 250 ticks across 8 seeds. The endpoints in
$(\text{Gini}, \Phi, \text{alive})$ space cluster into three regimes,
labelled by k-means with $k=3$ (silhouette score 0.55, with the
collapse-vs-extractive boundary the least sharp;
Table~\ref{tab:regimes}).

\begin{figure}[t]
\centering
\begin{subfigure}{0.5\linewidth}
\includegraphics[width=\linewidth]{phase3d_trajectories.png}
\end{subfigure}\hfill
\begin{subfigure}{0.47\linewidth}
\includegraphics[width=\linewidth]{phase3d_shadows.png}
\end{subfigure}
\caption{Phase portrait of thirteen scenarios in (Gini, $\Phi$, alive).
Three regimes are visible: viable autocatalytic ($\Phi > 0.4$, alive
$> 65$), extractive equilibrium ($\Phi \approx 0.01$, alive 40--75),
and collapse ($\Phi = 0$, alive $< 20$). State excluded from Gini and
alive count.}
\label{fig:phase}
\end{figure}

\begin{table}[t]
\centering
\small
\caption{Endpoint coordinates by scenario at $N=100$, 250 ticks, mean
across 8 seeds, State excluded from Gini.}
\label{tab:regimes}
\begin{tabular}{lrrrl}
\toprule
scenario & Gini & $\Phi$ & alive & regime \\
\midrule
socialist             & 0.30 & 0.880 & 99.0 & viable \\
nordic                & 0.30 & 0.701 & 89.9 & viable \\
inheritance\_break\_ubi & 0.57 & 0.667 & 85.9 & viable \\
maker\_economy        & 0.77 & 0.586 & 68.8 & viable \\
tax\_and\_spend       & 0.43 & 0.567 & 83.6 & viable \\
memory\_equalised     & 0.92 & 0.015 & 76.9 & extractive \\
creditless            & 0.87 & 0.013 & 55.2 & extractive \\
current\_world        & 0.89 & 0.011 & 63.1 & extractive \\
memory\_starved       & 0.87 & 0.011 & 55.5 & extractive \\
stevenson\_world      & 0.88 & 0.001 & 20.2 & collapse \\
feudal                & 0.78 & 0.000 & 15.2 & collapse \\
inheritance\_break    & 0.00 & 0.000 &  0.0 & collapse \\
\bottomrule
\end{tabular}
\end{table}

The viable regime, which sits at $\Phi \geq 0.4$ and alive $\geq 65$
on average, is reached by three distinct policy mixes: heavy state
spending (\texttt{socialist}, \texttt{nordic}), high inheritance
taxation combined with redistribution (\texttt{inheritance\_break\_ubi},
\texttt{tax\_and\_spend}), or simply higher maker productivity
(\texttt{maker\_economy}). The collapse regime is reached by
configurations that maximise extraction (\texttt{stevenson\_world},
\texttt{feudal}) or that remove inheritance without a safety net
(\texttt{inheritance\_break}); in the last case the model collapses
within 30 ticks (Figure~\ref{fig:survival}).

\begin{figure}[t]
\centering
\includegraphics[width=0.97\linewidth]{survival_scenarios.png}
\caption{Survival curves for eight illustrative scenarios at $N=100$,
6 seeds. \textbf{Left:} alive over time. \textbf{Right:} cumulative
excess mortality as fraction of initial population.}
\label{fig:survival}
\end{figure}

%==========================================================================
\section{Visualising the flow regime: the Sankey of per-tick accrual}
\label{sec:sankey}
%==========================================================================

The Sankey diagram (Figure~\ref{fig:sankey}) renders mean per-tick
wealth flows in steady state as a network with edge widths proportional
to flow magnitude. Three policy configurations are shown side by side
to make $\Phi$ legible as a diagram.

\begin{figure}[t]
\centering
\includegraphics[width=\linewidth]{sankey_flows.png}
\caption{Steady-state wealth flows for three policies at
$N=10{,}000$, 3 seeds, ticks 80--180. \textbf{Left
(\texttt{uk\_tory\_2010\_24}):} the rentier--lender deposit loop
dominates (44\% of total flow); productive wage flow
(maker+employer to worker) is 11\%; $\Phi \approx 0.05$.
\textbf{Middle (\texttt{zucman\_2pct\_spend}):} the rentier--lender
loop has shrunk; wealth-tax revenue is recirculated to workers via
state employment and purchases; $\Phi$ is around 0.44.
\textbf{Right (\texttt{socialist}):} worker--state--worker
circulation accounts for 40\% of total flow; the rentier--lender
loop has effectively disappeared; $\Phi \approx 0.88$.}
\label{fig:sankey}
\end{figure}

The Sankey is the cleanest single artefact for explaining $\Phi$ to
a non-technical audience. It also reads as a visualisation of the
productive-vs-extractive distinction independently of the specific
$\Phi$ definition.

%==========================================================================
\section{Cultural erosion under economic pressure}
\label{sec:culture}
%==========================================================================

This section presents the paper's principal alife-relevant
contributions. The model exhibits cultural erosion rather than
cumulative cultural evolution, a property that follows directly from
its closed K=5 skill space. We show the erosion under wealth
concentration, decompose the bandwidth response into the
non-novel and novel parameter axes, and present the class-by-class
skill-to-wealth elasticity decomposition that grounds the
cultural-economic asymmetry in mechanism.

\subsection{The memory sweep: breadth varies fourfold, Gini barely moves}

\begin{figure}[t]
\centering
\includegraphics[width=0.97\linewidth]{memory.png}
\caption{The \texttt{memory} sweep at $N=100$. We vary base bandwidth
$b_0 \in [0.2, 1.0]$ and the wealth-coupling slope $b_w \in [0, 1]$
on a $9 \times 9$ grid, 3 seeds. Cultural breadth (top middle) varies
by a factor of four, from 1.3 to 5.0. Gini (top left) varies by only
10\%, from 0.84 to 0.92.}
\label{fig:mem}
\end{figure}

We sweep the two bandwidth parameters in
Equation~\ref{eq:bandwidth}, $(b_0, b_w)$, over $[0.2, 1.0] \times
[0, 1]$. Across this range cultural breadth (mean count of nonzero
skill dimensions across agents) varies by a factor of four. Gini
(state excluded) varies by only $10$\%.
\texttt{memory\_equalised} (full bandwidth, no wealth coupling)
gives the \emph{highest} non-extractive Gini in the table; equalising
cultural transmission alone does not equalise wealth.

This is the cultural-economic asymmetry in its raw form. The
remainder of this section explains it.

\subsection{Bandwidth decomposition: the wealth coupling is a
real, secondary channel}
\label{sec:bandwidth-decomp}

\begin{figure}[t]
\centering
\includegraphics[width=0.97\linewidth]{bandwidth_decomp.png}
\caption{Decomposition of the \texttt{memory} sweep into the
$b_0$ and $b_w$ effects. \textbf{Left:} cultural breadth on the
$(b_0, b_w)$ plane. \textbf{Middle:} marginal effect of $b_0$
(averaged over $b_w$) and of $b_w$ (averaged over $b_0$).
The $b_0$ axis dominates with range 2.99 dimensions; the
wealth-coupling $b_w$ contributes a graded response with range
0.87 dimensions across the full wealth distribution. \textbf{Right:}
realised distribution of bandwidth\_dims at birth by class under
\texttt{current\_world}; higher-wealth classes routinely inherit
4--5 dimensions, workers cluster at 2--3.}
\label{fig:bandwidth}
\end{figure}

We decompose the breadth response into marginal effects of $b_0$
and $b_w$ (Figure~\ref{fig:bandwidth}, middle).
The $b_0$ axis dominates: averaging over $b_w$, varying $b_0$
across $[0.2, 1.0]$ moves cultural breadth by 2.99 dimensions
(from 2.0 to 5.0). The wealth-coupling $b_w$ axis contributes a
graded response of 0.87 dimensions (from 3.2 to 4.1 on average
across the $b_0$ range).

The wealth coupling is therefore a real, secondary channel
contributing roughly $0.87 / (0.87 + 2.99) = 22$\% of the breadth
variation around the base. The previous (linear) coupling formulation
saturated above parent wealth $\approx 500$ units (\pounds 50M
under the unit mapping used in the companion economic-policy paper)
and was effectively a step function on poor-vs-everyone-else; the
log-wealth coupling stays graded across orders of magnitude. The
realised distribution of bandwidth\_dims at birth
(Figure~\ref{fig:bandwidth}, right) shows the coupling actually
operating: workers cluster at 2--3 dimensions and rentiers and
employers routinely inherit 4--5.

\subsection{The asymmetry is mechanism, not architecture: skill-to-
wealth elasticity by class}
\label{sec:elasticity}

We test directly whether the model's cultural-economic asymmetry
is architectural. If the skill-to-wealth channel were uniformly
weak by construction across all classes, the negative result
(cultural transmission does not equalise wealth) would not be
informative.

We regress 40-tick wealth gain on the five-dimensional skill
vector, by class. At $N = 500$ (population\_scale = 5), 8 seeds,
WARMUP $= 20$ ticks dropped as transient, 40-tick window, we obtain
$n = 809$ agent-windows with adequate sample size in every class.
Table~\ref{tab:elasticity} reports normalised elasticities
$\eta_i = \beta_i \sigma_{\text{skill}_i} /
|\overline{\text{gain}}|$, $R^2$, and sample size $n$.

\begin{table}[t]
\centering
\small
\caption{Skill-to-wealth normalised elasticities by class. Wealth gain
over 40 ticks regressed on the five-dimensional skill vector at
$N=500$.}
\label{tab:elasticity}
\begin{tabular}{lrrrrrrl}
\toprule
class & tech & org & social & fin & gen & $R^2$ & $n$ \\
\midrule
WORKER     & 0.76 & 0.95 & $-0.02$ & $-0.03$ & $-0.03$ & 0.97 & 395 \\
MAKER      & 0.00 & 0.00 & 0.01 & 0.00 & 0.33 & 0.98 & 144 \\
EMPLOYER   & 0.30 & 0.03 & 0.25 & $-0.10$ & 0.24 & 0.12 & 100 \\
LENDER     & $-0.01$ & 0.00 & $-0.01$ & 0.00 & $-0.01$ & 0.13 & 79 \\
RENTIER    & $-0.06$ & $-0.05$ & 0.09 & $-0.10$ & $-0.08$ & 0.12 & 47 \\
SPECULATOR & 0.56 & $-0.39$ & $-1.41$ & 0.04 & 0.13 & 0.11 & 44 \\
\bottomrule
\end{tabular}
\end{table}

The result is decisive. Workers have $R^2 = 0.97$ with
class-typical elasticities $\eta_\text{tech} = 0.76$ and
$\eta_\text{org} = 0.95$. Makers have $R^2 = 0.98$ with
$\eta_\text{gen} = 0.33$. Employers have moderate $R^2 = 0.12$.
Lenders, rentiers, and speculators have $R^2 \in [0.11, 0.13]$ and
$|\eta_i| \leq 0.10$ for lenders and rentiers.

The cultural-economic asymmetry is therefore not a wiring artefact.
One half of the economy has strong skill-to-wealth feedback and the
other half does not. The cultural transmission mechanism is doing
its share of work in the half of the economy where skills matter for
wealth (workers and makers), and the wealth dynamics that produce the
overall Gini are anchored in the half of the economy (rentiers and
lenders) where skills do not. The architecture builds in the bias of
the coupling (extractive classes earn from stocks rather than
skills), and the model surfaces the consequence as the observed
cultural-economic asymmetry.

This is the contribution we believe makes the model useful for
ALIFE. The negative result ``cultural transmission alone does not
equalise wealth'' is well established in the social-science
literature in various forms; what we add is a mechanism-level
account of why, in a stylised ABM, the negative result is robust:
because the wealth dynamics of the dominant extractive classes are
independent of their skills.

\subsection{Direct test of \citet{spangler2022distribution}}
\label{sec:spangler}

The closest published ABM neighbour finds that skill inheritance
raises wealth inequality. We can test that mechanism in our model
under a clean toggle that holds bandwidth-gating and class-bias
persistence intact and swaps only whether transmitted skill values
are parent-derived or drawn from the class-typical prior. A coarser
toggle, which destroys the bandwidth mechanism along with parental
transmission by setting the mutation rate $\mu = 1.0$, is reported
for contrast.

At $N = 100$, 10 seeds, 250 ticks:

\begin{table}[h]
\centering
\small
\caption{Spangler \& Sarkar test: clean toggle vs.\ confounded
toggle. Medians.}
\label{tab:spangler}
\begin{tabular}{lrrr}
\toprule
metric & default & class-prior & no-inheritance \\
       & (parent-derived) & (clean toggle) & ($\mu = 1$, confounded) \\
\midrule
Gini             & 0.885 & 0.876 & 0.908 \\
Top-1 share      & 0.235 & 0.250 & 0.254 \\
Top-10 share     & 0.785 & 0.683 & 0.891 \\
Alive            & 61.5  & 56.5  & 73.0 \\
Cultural breadth & 3.38  & 3.55  & 5.00 \\
\bottomrule
\end{tabular}
\end{table}

Under the clean toggle, swapping parent-derived for class-typical-
prior transmitted values \emph{lowers} Gini by 0.010 and \emph{raises}
top-1 share by 0.016. So parent-derived skill inheritance does raise
Gini relative to a class-typical prior; the direction is a partial
reproduction of \citet{spangler2022distribution}, with a magnitude
much smaller than they report.

The mechanism for the small magnitude is in
Table~\ref{tab:elasticity}: skill inheritance acts strongly on
workers and makers but those classes are the ones whose wealth is
already most closely tied to skills, so swapping the transmitted
value for a class-typical prior produces only a moderate
within-class shift. The wealth dynamics of the rentier and lender
classes, which dominate the overall Gini, are insensitive to skill
content. A single-channel ABM in which skills feed wealth via a
uniform mechanism (Spangler \& Sarkar's design) is expected to
amplify the inheritance effect.

%==========================================================================
\section{Open-ended skill space: a minimal demonstration}
\label{sec:open-skills}
%==========================================================================

The cultural section above is honest that the K=5 skill space is
closed and that the model produces erosion, not cumulative
evolution. For an ALIFE audience this is uncomfortable framing on
its own, and we add a stand-alone open-ended-evolution demonstration
that addresses the open-endedness concern at demonstration scale.

In a separate stand-alone simulation the skill vector is variable-
length. At every reproduction event, with probability proportional
to $\log(w_\text{parent})$, a previously-undiscovered skill
dimension enters the global pool and is initialised on the child to
a fresh Beta(2,2) value. Subsequent inheritance follows the same
bandwidth rules as the main model. The wealth-coupling of novelty
events is the same coupling that the bandwidth gate operationalises.
500 initial agents, 400 ticks, 5 seeds.

\begin{figure}[t]
\centering
\includegraphics[width=0.97\linewidth]{open_skills.png}
\caption{Open-ended skill-space demonstration. \textbf{Left:}
active skill-pool size over time. The closed variant (grey, no
novelty) stays at five; the open variant (blue, novelty proportional
to log parent wealth) climbs from five to about ten over 400 ticks.
\textbf{Middle:} per-agent cultural breadth grows accordingly.
\textbf{Right:} per-capita discovery rate by class
(discoveries per agent-generation).}
\label{fig:open-skills}
\end{figure}

Two results.

First, the skill space grows under wealth-biased novelty: the
active pool climbs from $5$ to about $10$ over 400 ticks under the
open variant. The closed variant stays at $5$ by construction. The
mechanism can therefore be extended to genuine open-ended novelty
within its existing architecture.

Second, on a per-capita basis the wealth coupling does concentrate
discovery in wealthier lineages. With approximate average class
populations of 140 workers, 40 makers, and 20 rentiers per run, and
50 discoveries observed across 5 seeds, the per-capita discovery
rates (discoveries per agent-generation, where generation length is
estimated as median lifespan / ticks) are 0.0092 for workers,
0.0052 for makers, and 0.0123 for rentiers. Rentiers discover at
$1.3 \times$ the worker rate and $2.4 \times$ the maker rate on a
per-capita basis. Raw counts are dominated by worker abundance and
do not support a disproportionate-contribution claim on their own;
the per-capita rate does.

We do not claim this as a load-bearing finding. It is a
demonstration scoped to address the concern that the main model's
closed skill space is a category problem for an ALIFE audience. A
more substantive treatment that folds the open-ended discovery
mechanism back into the full economic model is the natural follow-up.

%==========================================================================
\section{Robustness and extensions}
\label{sec:robustness}

We have done three additional checks since the first version of
this work.  Each addresses an open-endedness or robustness concern
that the cultural-evolution audience commonly raises.

\subsection{Elasticity is robust across scenarios}
\label{sec:elasticity-robustness}

The skill-to-wealth elasticity decomposition
(Section~\ref{sec:elasticity}, Table~\ref{tab:elasticity}) was
reported on the \texttt{current\_world} scenario.  The class-
specific pattern that grounds the cultural-economic asymmetry
might be specific to that scenario; we test by rerunning the
elasticity regression at \texttt{nordic},
\texttt{maker\_economy}, and \texttt{feudal}.

\begin{table}[h]
\centering
\small
\caption{Skill-to-wealth $R^2$ by scenario and class at $N=500$,
8 seeds.  Patterns are qualitatively unchanged across scenarios.}
\label{tab:elasticity-scenarios}
\begin{tabular}{llrrrrrr}
\toprule
scenario & metric & WORKER & MAKER & EMPLOYER & LENDER & RENTIER & SPECULATOR \\
\midrule
current\_world & $R^2$ & 0.97 & 0.98 & 0.12 & 0.13 & 0.12 & 0.11 \\
nordic         & $R^2$ & 0.71 & 0.23 & 0.08 & 0.12 & 0.11 & 0.13 \\
maker\_economy & $R^2$ & 0.91 & 0.98 & 0.04 & 0.16 & 0.06 & 0.05 \\
feudal         & $R^2$ & 0.71 & 0.47 & 0.02 & 0.28 & 0.02 & 0.01 \\
\bottomrule
\end{tabular}
\end{table}

The qualitative pattern holds: workers and makers have strong
skill-to-wealth fits; employers, lenders, rentiers, and
speculators have weak ones.  The maker $R^2$ drops to $0.23$ in
\texttt{nordic} (where the state-employment subsidy reduces the
within-maker variance attributable to skill) and to $0.47$ in
\texttt{feudal} (where extraction dominates).  The asymmetry
between productive and extractive classes survives intact.

\subsection{Long-run cumulative novelty}
\label{sec:longrun-novelty}

The open-skills demonstration in Section~\ref{sec:open-skills}
ran for 400 ticks.  An open-ended-evolution claim requires that
discovery be sustained beyond the initial transient.  We rerun the
demonstration at horizons $400$, $800$, and $1600$ ticks with
$K_\text{max} = 25$.

\begin{table}[h]
\centering
\small
\caption{Long-run open-skill discovery at three horizons, 5 seeds.}
\label{tab:longrun-novelty}
\begin{tabular}{rrrr}
\toprule
horizon & final K (mean) & final K (std) & dims discovered per 100 ticks \\
\midrule
400  & 16.00 & 4.34 & 2.75 \\
800  & 25.00 & 0.00 & 2.50 \\
1600 & 25.00 & 0.00 & 1.25 \\
\bottomrule
\end{tabular}
\end{table}

The active skill pool reaches the $K_\text{max} = 25$ cap by
horizon $800$ at every seed, and the per-100-tick discovery rate
remains above $1$ at horizon $1600$ (limited by the cap, not by
the underlying mechanism).  The demonstration is therefore not
one-shot near-equilibrium; the discovery mechanism is sustained.
A truly cap-free OEE claim would require raising $K_\text{max}$,
which the framework permits but which we do not pursue here.

\subsection{Hard material constraint: an autopoietic reading of $\Phi$}
\label{sec:material}

The model so far treats material inputs to makers as an unlimited
constant; the maker pays a material cost per tick, but the cost
is drawn from no real reservoir.  To test the autopoietic reading
of $\Phi$ (i.e.\ that high $\Phi$ characterises a system
sustaining its own input flows), we add a finite global material
pool that depletes via maker production and replenishes at a
fixed rate per tick.

At regen rate $0$, the pool of $50$ units exhausts within roughly
$25$ ticks and the maker class can no longer produce; alive at
$N=100$ falls to $40$ from baseline $64$.  At regen rate $5$
units per tick, alive holds at $54$ with $\Phi$ at $0.026$.  At
regen rate $20$ or above, alive reaches its full baseline of
$64$--$66$ with $\Phi$ at $0.032$.  The same regen sweep on
\texttt{maker\_economy} gives a steeper coupling: alive $47, 54,
72, 72$ across regen $\{0, 5, 20, 100\}$, $\Phi$ $0.519, 0.639,
0.491, 0.491$.

The finding is that $\Phi$ at sustained-flow values requires a
material supply that the production sector can keep ahead of.
The autopoietic reading is supported: high $\Phi$ is a property of
a system that is materially viable as well as monetarily so.  The
model can fail by running out of material, not only by running out
of wealth.

This is a step toward the genuine alife reading of the closure
metric.  The result is preliminary; the material model is a single
global pool rather than a spatially or class-locally structured
reservoir, and the regen mechanism is exogenous.  A more
substantive treatment with class-local material reservoirs is the
natural next step.

\section{Discussion}
\label{sec:discussion}
%==========================================================================

\subsection{What we believe is novel for an ALIFE audience}

The class-specific skill-to-wealth elasticity decomposition
(\S\ref{sec:elasticity}, Table~\ref{tab:elasticity}) is the freshest
contribution. The cultural-economic asymmetry the bandwidth
coupling generates is not a structural artefact of how the model is
wired; it is a consequence of one half of the economy having near-
zero skill-to-wealth feedback while the other half has strong
feedback. We believe this kind of mechanism-level grounding of a
negative result is exactly the type of result ALIFE values.

The wealth-coupled cultural-transmission bandwidth mechanism
(\S\ref{sec:transmission}) is, to our reading of the cultural-
transmission and computational-cultural-evolution literatures, not
present in any published ABM. \citet{henrich2001evolution} couples
the \emph{target} of transmission to status; population-size
constraints \citep{henrich2004demography, lewis2012transmission}
discuss aggregate fidelity. The bandwidth-coupled-to-wealth coupling
here appears to be genuinely new.

The productive-flow share $\Phi$ (\S\ref{sec:phi}) is a continuous,
scalar, per-tick measurement device for flow regime. We are explicit
that it is not a topological autocatalytic-set closure
\citep{kauffman1988evolution, hordijk2004detecting}; it is a
flow-share that operationalises the productive-vs-extractive
distinction.

The open-skills demonstration (\S\ref{sec:open-skills}) shows that
the same wealth coupling that gates transmission can be extended to
seed novelty, and that on a per-capita basis the coupling
concentrates discovery in wealthier lineages.

\subsection{What the model cannot do}

The K=5 skill space is closed in the main results. The open-ended
extension is stand-alone and does not currently fold back into the
full economic model. The economic layer is itself non-open-ended in
ways that constrain the cultural reading: agent-type proportions are
fixed, the employer count does not change, no new occupational
categories can emerge. For an ALIFE audience these are real
limitations that we acknowledge.

The K=5 grain choice for skills is a methodological commitment
following \citet{charbonneau2021fidelity}'s sense of the
``grain problem'' for fidelity measures: the dimensions are not
natural primitives but a chosen discretisation. The model's results
should be expected to be robust to grain refinement, and that
robustness is not tested here.

The skill-to-wealth elasticity is measured under a single scenario
(\texttt{current\_world}) and a single time window (40 ticks). The
class-specific structure is qualitatively robust across the
\texttt{maker\_economy} and \texttt{tax\_and\_spend} variants we
have spot-checked, but a full sensitivity analysis is future work.

\subsection{Honest reading}

The model is presented as an intuition pump for cultural-economic
dynamics, not a predictive instrument. The K=5 skill space is small
enough to read but large enough that the wealth-coupled bandwidth
mechanism has a non-trivial decomposition. The closed-space limit is
acknowledged. The open-skills extension shows where the next paper
will land.

The cultural-economic asymmetry as we have characterised it is the
asymmetry between productive and extractive classes in their
internal skill-to-wealth feedback. The model says, and the
elasticity table supports, that for an audience interested in the
emergence of inequality from coupled cultural and economic
mechanisms, the architecture of the economic layer matters as much
as the architecture of the cultural layer. The cultural channel,
in this stylised setting, is the lesser of the two contributors to
the overall outcome, and is dominated by the wealth-feedback
asymmetry across the productive/extractive split.

%==========================================================================
\section{Conclusion}
%==========================================================================

We have presented an agent-based model in which a wealth-coupled
bandwidth gate determines cultural transmission, and a continuous
flow-share metric $\Phi$ characterises whether per-tick monetary
accrual is dominated by productive or extractive activity. The
cultural side of the model exhibits erosion, not cumulative
evolution; a stand-alone open-ended extension demonstrates the
mechanism's natural direction for further development. The
cultural-economic asymmetry the model produces is grounded in
mechanism via a class-specific skill-to-wealth elasticity
decomposition. The model is honest about scope: a stylised
intuition pump for coupled cultural-economic dynamics, not a
predictive instrument.

The fiscal-policy implications of the model (wealth-tax design,
inheritance dynamics, calibration to UK wealth-share data) are
reported separately. The contribution this paper foregrounds is
the cultural mechanism and the elasticity decomposition that
grounds its emergent asymmetry.

\bibliographystyle{plainnat}
\bibliography{references}

\appendix
\section{ODD-style summary of the model}
\label{app:odd}

A brief ODD (Overview, Design concepts, Details) summary. A full
ODD+D specification and the companion economic-policy paper are
deferred to a separate supplement.

\paragraph{Overview.}
Seven agent classes (WORKER, MAKER, EMPLOYER, LENDER, RENTIER,
SPECULATOR, STATE) at fixed proportions per 100 agents. Each agent
has wealth, debt, age, lifespan, a five-dim skill vector, alive flag,
peak lifetime wealth, bandwidth-at-birth.

\paragraph{Process.}
Per tick: production (makers, employers, optional state
employment); extraction (rent, capital return, speculator returns);
consumption and credit; interest and repayment; lender deposits and
optional state recapitalisation; tax and UBI; life (age, death,
reproduction with bandwidth-limited inheritance).

\paragraph{Cultural mechanism.}
At reproduction, $\text{bw\_dims}(w_{\text{parent}}) = \text{clip}(
b_0 K + b_w K \log_{10}(\max(1, w_{\text{parent}}))/4, 1, K)$ with
$K=5, b_0=0.5, b_w=0.5$. Strongest dimensions transmitted;
per-dim mutation rate 0.05.

\paragraph{Metrics.}
Per tick: Gini (state-excluded by default), $\Phi$, alive, cultural
breadth, cultural diversity (variance), knowledge lost, top-share
percentiles (raw and mortality-corrected).

\paragraph{Code.} \texttt{sim.py} for the core model.
\texttt{scenarios.py} for the 13-scenario library. \texttt{elasticity.py}
for the regression. \texttt{open\_skills.py} for the open-ended
demonstration. Default seed 42, default population\_scale 1.
\end{document}
